\section{Resultados}

En esta sección se presentan los hallazgos cuantitativos y cualitativos obtenidos. Siguiendo las directrices de APA 7, los resultados estadísticos se reportan con la notación formal pertinente y las tablas se integran sin líneas verticales divisorias.

\subsection{Análisis Descriptivo}
En la Tabla~\ref{tab:rendimiento} se resumen las métricas de participación y desempeño registradas en ambos grupos experimentales.

\begin{table}[h]
    \caption{Resultados Obtenidos en las Pruebas de Rendimiento por Grupo}
    \label{tab:rendimiento}
    \centering
    \begin{tabular}{lccc}
        \toprule
        Grupo & Muestra ($n$) & Media ($M$) & Desviación Estándar ($DE$) \\
        \midrule
        Control & 50 & 74.20 & 8.45 \\
        Experimental & 50 & 86.60 & 6.12 \\
        \bottomrule
    \end{tabular}
    \begin{tablenotes}
        \small
        \item \textit{Nota.} Los puntajes corresponden a la escala estandarizada de evaluación (0 a 100 puntos).
    \end{tablenotes}
\end{table}

\subsection{Pruebas de Hipótesis y Significancia}
Los análisis comparativos reflejaron diferencias estadísticamente significativas entre el grupo experimental y el grupo de control, $t(98) = 8.41$, $p < .001$, con un tamaño del efecto considerable, $d = 1.68$. Estos valores respaldan la eficacia del esquema propuesto.
